\chapter{Accessibility in the terminal}
\label{ch:a11y2}

\begin{aside}
Accessibility in TUIs is harder than GUIs because screen
readers don't see ANSI escapes. The mitigations are: keyboard-
only navigation, semantic labels, high contrast, and avoiding
flicker.
\end{aside}

\section{Screen reader compatibility}

Most TUI screen readers (BRLTTY, Speakup) treat the screen
as plain text. They can't follow ANSI movement; they read
the current row when the cursor moves.

The implication: organise your UI so that reading rows
top-to-bottom makes sense. Don't rely on visual proximity
across columns.

\section{Keyboard-only}

Every action must have a keyboard path. Mouse is a
convenience; not a requirement. \maya{} enforces this for
its built-in widgets; custom widgets must follow.

\section{Focus order}

Tab cycles through focusable elements. The order should
follow reading order (top-to-bottom, left-to-right) by
default. Override only when content order disagrees with
intent.

\section{Labels}

A text input with a visible label next to it is fine for
sighted users. For screen readers, the relationship is
implicit. Make it explicit:

\begin{lstlisting}
auto field = text_input(value)
           | accessible_label("Email");
\end{lstlisting}

The label is reported when focus enters.

\section{Live regions}

Updates that the user needs to know about (notifications,
status changes) should be announced. \maya{}'s
\code{live\_region} marks an element so screen readers
announce its content changes.

\section{Reduced motion}

Respect \code{PREFERS\_REDUCED\_MOTION} or an in-app
setting. Disable animations.

\section{High contrast}

Provide a high-contrast theme. The default themes meet AA
contrast (4.5:1); a dedicated AAA theme is a useful add.

\section{Avoiding flicker}

Sub-frame artefacts (blink-off, mid-frame redraw) trigger
seizures in some users. Use synchronised output (DEC 2026)
when available. Avoid rapid colour changes.

\section{NO\_COLOR}

Respect \code{NO\_COLOR} for users with colour vision
deficiency or screen-reader compatibility.

\section{Recap}

\begin{recap}
\begin{itemize}[leftmargin=*]
  \item Reading order = logical order.
  \item Every action has a keyboard path.
  \item Labels explicit, not implicit.
  \item Live regions announce updates.
  \item Reduced motion, high contrast, \code{NO\_COLOR}.
\end{itemize}
\end{recap}

\section*{Workshop}

\begin{enumerate}[leftmargin=*]
  \item Audit your app: every action keyboard-reachable?
  \item Add accessible labels to every input.
  \item Mark a status area as a live region.
\end{enumerate}
